\section{Discussion}\label{sec:discussion}
\finding{Residuals find the model's margin, including the part that is not theirs to spend.} The unshielded policy did
exactly what its reward asked: it found that the friction model over-predicts the grip needed for a fast lift and spent
that margin. Nothing in its training signal distinguished the margin that comes from quasi-static conservatism, which is
real, from the margin that protects against a friction drop, which is needed only when the drop occurs. With $r$ drawn from
$[1,1.2]$ during training, the cut paid off on most batches and failed on a few (Fig.~\ref{fig:training}). We expect the
same pattern wherever a learned correction sits on a conservative model of a contact transition: the conservatism is
visible in the average reward, the transition only in the tail.

\finding{Certify onset, not slip.} The shield removed every failure and cost 6\% grip. But it did so by forbidding slip,
not by bounding it: every certified lift had a bound of zero. Section~\ref{sec:q4} and Appendix~\ref{app:cliff} show why
this is not an accident of our tolerances. With any friction drop, the kinetic deficit outlasts the moment of onset, so
the certified slip jumps from zero to centimetres. For lifting, a useful certificate is therefore one about \emph{onset}:
$g+a(t)\le gF/\Fs$ throughout the profile. It is simple, it needs no $r_{\max}$ at all, and our results suggest it gives up
almost nothing relative to the full slip bound once $r_{\max}\ge1.1$. A slip allowance becomes valuable only when the
friction drop is known to be small, for example on materials whose static and kinetic coefficients are measured to be
close. The extended study adds one qualification: a certificate that sits exactly on the onset boundary inherits whatever
happens there, and when friction weakened during pre-sliding creep (row E9) the certified lifts failed. An onset
certificate should carry a margin, sized from measurements of how a contact creeps near its static limit.

\finding{Grip or time.} When the shield must act, raising the grip kept the speed and slowing the lift kept the grip. Which
is better depends on the application: for fragile objects the time projection is the natural choice, for throughput the
grip projection. Both remove failures; neither declined in our study.

\finding{Where the time goes.} The lift is a small part of an episode. Halving the probe time would save more than the
learned lift does. Probes that use the lift itself, for example a slow initial acceleration that doubles as a breakaway test,
are a promising direction.

\section{Limitations}\label{sec:lim}
\finding{Simulation only.} All results are from MuJoCo with an emulated stick--slip transition. Real contacts show
rate-dependent friction, pre-sliding displacement, load-dependent soft-finger friction and wear~\cite{canudas1995,xydas1999};
our switching thresholds were tuned in development, and the switch reads the relative speed with a one-step lag that we
found only during the extended study. The lag did not change any confirmatory result (row R0L) but caused the
certificate failures of row E7. The certificate assumes that $\Fs$ bounds the breakaway grip during the
lift, which held here because the quasi-static model was conservative for fast lifts (Table~\ref{tab:val}); on hardware this
must be measured. Domain randomisation~\cite{tobin2017,peng2018} over friction models would be a natural first step before
hardware.
\finding{A probe with a blind spot.} The probe assumes $\mu_s\ge0.3$ and a slip monitor much quieter than its motion
threshold; when either fails, it reports the maximum grip and the shield declines (rows E2, E3). Safe, but not gentle.
\finding{Narrow task.} The rod is a rigid box constrained to slide axially, held at its centre of mass by two pads. Rotational
slip, off-centre grasps, deformable objects and multi-finger hands are outside this study.
\finding{Policy class and budget.} Policies are linear and were trained for 1{,}560 episodes; all of them saturated the duration
action. Training choices matter: twice the budget made the unshielded policy nearly safe on our test set, while training
without a friction drop made it much less safe (Section~\ref{sec:trainrows}). Richer policies might find other trade-offs.
\finding{Deviations.} The pre-registered exceedance metric counts elastic creep (D1); the confirmation set was analysed
pooled (D2); the slip-allowance analysis (D3) and the whole extended study are exploratory. Appendix~\ref{app:protocol} lists
them.

\section{Conclusion}
We asked whether a learned lifting policy should be allowed to trade a little slip for a lighter, faster lift. A residual
policy on a friction-model controller lifted 2.7 times faster; left unshielded, it also learned to spend a safety margin it
could not see, and failed whenever friction dropped. A closed-form slip bound used as a shield removed every failure in
simulation at 6\% extra grip and no loss of speed, even beyond its assumed friction drop, and the results replicated on
fresh objects. The slip allowance itself turned out to be worth little: once friction may drop by 10\% or more, the only
certifiable lifts are those in which slip never starts. Our recommendation for practice is to certify onset and to spend
learning effort on faster probing. The next step is to measure the probe, the bound and the friction drop on hardware.

\section*{Broader Impact}
This work aims to make robot grasping gentler and safer. Its main risk is over-trust: a certificate proved in simulation, or
under an assumed friction drop, can fail on real objects, as our own $r_{\max}=1.0$ results show. We report the conditions
under which the certificate held and failed so that it is not applied beyond them.

\section*{Data Availability}
Code, the frozen protocol with its addendum and deviations, trained weights, per-episode results and the scripts behind every
number and figure: \url{https://github.com/CAN-Aidoo/lift-shield}, archived at
\href{https://doi.org/10.5281/zenodo.23769371}{doi:10.5281/zenodo.23769371}.

\section*{Use of AI Tools}
Code, analysis and drafting were carried out with the assistance of an AI model (Claude, Anthropic) under the author's
direction. The author takes responsibility for the content.
